% ============================================================ %
%  Activity V: Public Key Infrastructure                        %
%  Styling lives in preamble.tex, per-document fields live in  %
%  metadata.tex. This file is the actual document content.     %
% ============================================================ %

% ============================================================ %
%  PREAMBLE                                                     %
%  Packages, fonts, spacing, and reusable macros only.          %
%  Per-document content lives in metadata.tex and template.tex. %
% ============================================================ %

% DOCUMENT CLASS SETUP %
\documentclass[12pt, a4paper]{article}

% DOCUMENT PACKAGES %
\usepackage[a4paper, margin=1in]{geometry}
\usepackage{enumitem}
\usepackage{amsmath}
\usepackage{amssymb}
\usepackage{mathtools}
\usepackage{graphicx}
\usepackage{xcolor}
\usepackage{listings}
\usepackage{hyperref}
\usepackage[capitalize, noabbrev]{cleveref}
\usepackage{csquotes}
\usepackage{tocloft}
\usepackage{array}
\usepackage{booktabs}
\usepackage{caption}
\usepackage{subcaption}
\usepackage{mwe}                 % provides \includegraphics{example-image*} placeholders
\usepackage{fancyhdr}
\usepackage{etoolbox}            % powers the optional bibliography toggle

% READABILITY / "AIRY" TYPOGRAPHY SETUP %
\usepackage{setspace}
\setstretch{1.15}
\setlength{\parindent}{0pt}
\setlength{\parskip}{0.6em plus 0.2em}
\usepackage{titlesec}
\titlespacing*{\section}{0pt}{2.4em}{1em}
\titlespacing*{\subsection}{0pt}{1.8em}{0.8em}
\titlespacing*{\subsubsection}{0pt}{1.4em}{0.6em}
\setlist{itemsep=0.4em, topsep=0.6em, parsep=0pt}

% FONT AND ENCODING SETUP %
\usepackage{fontspec}
\usepackage{polyglossia}
\usepackage[Thai]{ucharclasses}
\setdefaultlanguage{english}
\setotherlanguage{thai}
\setTransitionsFor{Thai}{\thaifont}{\normalfont}
\newfontfamily{\thaifont}{THSarabun.ttf}[
    Path=./fonts/,
    BoldFont=THSarabunBold.ttf,
    ItalicFont=THSarabunItalic.ttf,
    BoldItalicFont=THSarabunBoldItalic.ttf,
    Scale=1.23
]
\XeTeXlinebreaklocale "th"
\XeTeXlinebreakskip = 0pt plus 1pt

% Bibliography is a manual thebibliography in main.tex so Tectonic's
% bundled biblatex/biber versions do not have to match.

% CODE BLOCK PACKAGES %
% minted needs Pygments and shell escape, so build with:
%   tectonic -X compile -Z shell-escape main.tex
\usepackage{tcolorbox}
\tcbuselibrary{minted, skins, breakable}

% DEFINE GITHUB DARK THEME COLORS %
\definecolor{ghback}{HTML}{0d1117}
\definecolor{ghtext}{HTML}{c9d1d9}
\definecolor{ghborder}{HTML}{30363d}

% CODE BLOCK CONFIGURATION %
\renewcommand{\theFancyVerbLine}{\textcolor{white!60!gray}{\scriptsize\arabic{FancyVerbLine}}}
\newtcblisting{codeblock}[3][]{
    enhanced,
    breakable,
    colback=ghback,
    colframe=ghborder,
    coltitle=ghtext,
    colupper=ghtext,
    fonttitle=\bfseries\small,
    arc=3mm,
    boxrule=0.5mm,
    % PADDING SETTINGS %
    left=25pt,
    right=10pt,
    top=5pt,
    bottom=5pt,
    % MINTED SETTINGS %
    listing engine=minted,
    listing only,
    minted language=#2,
    minted options={
        linenos,
        numbersep=5pt,
        fontsize=\footnotesize,
        breaklines=true,
        breakanywhere=true,
        % autogobble drops the shared source indentation, so code can sit
        % nested inside the surrounding LaTeX; stripnl then keeps the
        % indented \end tag from adding a trailing blank line.
        autogobble,
        stripnl,
        style=monokai,
    },
    title={#3},
    #1
}

% SEMANTIC CALLOUT BOXES %
% Kept to a small, muted set on purpose: Note, Definition, Example, Takeaway.
% Each takes an optional argument to override the default title text, e.g.
%   \begin{notebox}[Note: edge case] ... \end{notebox}
\newtcolorbox{notebox}[1][Note]{
    enhanced, breakable,
    colback=blue!4, colframe=blue!45!black,
    colbacktitle=blue!45!black, coltitle=white,
    fonttitle=\bfseries, title=#1,
    arc=1.5mm, boxrule=0.4mm,
    left=8pt, right=8pt, top=5pt, bottom=5pt,
}
\newtcolorbox{defbox}[1][Definition]{
    enhanced, breakable,
    colback=black!3, colframe=black!55,
    colbacktitle=black!55, coltitle=white,
    fonttitle=\bfseries, title=#1,
    arc=1.5mm, boxrule=0.4mm,
    left=8pt, right=8pt, top=5pt, bottom=5pt,
}
\newtcolorbox{examplebox}[1][Example]{
    enhanced, breakable,
    colback=green!4, colframe=green!30!black,
    colbacktitle=green!30!black, coltitle=white,
    fonttitle=\bfseries, title=#1,
    arc=1.5mm, boxrule=0.4mm,
    left=8pt, right=8pt, top=5pt, bottom=5pt,
}
\newtcolorbox{takeawaybox}[1][Key Takeaway]{
    enhanced, breakable,
    colback=orange!5, colframe=orange!55!black,
    colbacktitle=orange!55!black, coltitle=white,
    fonttitle=\bfseries, title=#1,
    arc=1.5mm, boxrule=0.4mm,
    left=8pt, right=8pt, top=5pt, bottom=5pt,
}

% DISABLE PAGE NUMBER IN TABLE OF CONTENTS %
\tocloftpagestyle{empty}

% RUNNING HEADER / FOOTER (body pages only; cover and TOC stay empty) %
% Left = course code, center = current section, right = author.
\pagestyle{fancy}
\setlength{\headheight}{14pt}
\fancyhf{}
\fancyhead[L]{\small\scshape \coursecode \ \coursename}
% \fancyhead[C]{\small\itshape \rightmark}
\fancyhead[R]{\small \docauthor}
\fancyfoot[C]{\small\thepage}
\renewcommand{\headrulewidth}{0.4pt}
\renewcommand{\footrulewidth}{0pt}
\fancypagestyle{plain}{
    \fancyhf{}
    \fancyfoot[C]{\small\thepage}
    \renewcommand{\headrulewidth}{0pt}
}

% ============================================================ %
%  METADATA                                                     %
%  Edit this file per document. Everything here is the         %
%  "content" you change; the styling stays in preamble.tex.    %
% ============================================================ %

% SHOW BIBLIOGRAPHY? %
% Set to false to hide the "References" section at the end of the document.
\newtoggle{showbibliography}
\toggletrue{showbibliography}

% COVER AND HEADER FIELDS %
\newcommand{\doctitle}{2110413 COMPUTER SECURITY}
\newcommand{\docauthor}{Worralop Srichainont}
\newcommand{\docdate}{\today}
\newcommand{\coursecode}{2110413}
\newcommand{\coursename}{COMPUTER SECURITY}
\newcommand{\assignmentname}{Public Key Infrastructure}
\newcommand{\institution}{Computer Engineering, Chulalongkorn University}
\newcommand{\assignmentnumber}{5}
\newcommand{\instructor}{Krerk Piromsopa, Ph.D.}
\newcommand{\studentid}{6632200221}
\newcommand{\duedate}{\today}

\usepackage{tikz}
\usetikzlibrary{arrows.meta, positioning}

% PDF METADATA AND LINK STYLING (needs the fields from metadata.tex) %
\definecolor{linkaccent}{HTML}{1a56db}
\hypersetup{
    pdftitle={\doctitle},
    pdfauthor={\docauthor},
    pdfsubject={\coursecode\ -- \coursename},
    colorlinks=true,
    linkcolor=linkaccent,
    citecolor=linkaccent,
    urlcolor=linkaccent,
    bookmarksnumbered=true,
}

\renewcommand{\subsectionmark}[1]{\markright{#1}}

\begin{document}

    \begin{titlepage}
        \centering
        \thispagestyle{empty}
        \vspace*{4em}

        {\scshape Assignment}

        \vspace{3em}
        {\Huge\bfseries \doctitle}
        \vspace{1.5em}

        \noindent\rule{0.6\textwidth}{0.5pt}
        \vspace{1.5em}

        {\large Assignment \assignmentnumber} \\[0.4em]
        {\large \assignmentname}

        \vfill

        {\large \institution}
        \vspace{1em}

        {\large \docauthor}

        \vspace{0.4em}
        {\normalsize Student ID: \studentid}

        \vspace{2em}
        {\normalsize \docdate}
    \end{titlepage}

    \tableofcontents
    \thispagestyle{empty}
    \clearpage

    \pagenumbering{arabic}

    \section{Setup}

        Commands were run on Windows~11 with OpenSSL~3.5.9 (Miniconda), Python~3 with
        PyOpenSSL~26.4.0 and \texttt{pem}, and MiKTeX \texttt{xelatex}. The trusted root
        bundle \texttt{ca-certificates.crt} is the Mozilla CA list from
        \texttt{curl.se/ca/cacert.pem}. Scripts and raw outputs are in
        \texttt{code/pki/} and \texttt{code/pki/out/}. This OpenSSL has no default
        trust store, so the effect of \texttt{-CAfile} is visible directly. The observed
        certificates are the ones served on the day of the test and may change later.

    \section{Exercises 1--9: TLS connection and certificates}

        \subsection{Exercise 1: With and without \texttt{-CAfile}}

            \begin{codeblock}[]{bash}{Commands}
                openssl s_client -connect twitter.com:443
                openssl s_client -connect twitter.com:443 -CAfile ca-certificates.crt
                openssl s_client -connect twitter.com:443 -CApath empty_dir
            \end{codeblock}

            Typing \texttt{GET / HTTP/1.0} and Enter twice returned \texttt{HTTP/1.1 400 Bad Request}
            in all three runs. This is normal, since no valid host or resource was requested.

            \begin{table}[h]
                \centering
                \small
                \begin{tabular}{p{0.3\textwidth}p{0.36\textwidth}p{0.24\textwidth}}
                    \toprule
                    Command & Verify return code & Result \\
                    \midrule
                    no option & 20 (unable to get local issuer certificate) & \texttt{verify error} \\
                    \texttt{-CAfile ca-certificates.crt} & 0 (ok) & \texttt{Verification: OK} \\
                    \texttt{-CApath empty\_dir} & 20 (unable to get local issuer certificate) & \texttt{verify error} \\
                    \bottomrule
                \end{tabular}
                \caption{Result of the three \texttt{s\_client} runs.}
            \end{table}

            \textbf{Difference.} Both commands open the same encrypted TLS session (TLS~1.3,
            \texttt{TLS\_AES\_256\_\allowbreak GCM\_SHA384}) and receive the same certificate chain. The
            only difference is whether the client can \emph{authenticate} the server. Without
            \texttt{-CAfile} the client has no trusted root, so it cannot validate the chain and
            prints a \texttt{verify error}. With \texttt{-CAfile} it finds a trusted root and prints
            \texttt{Verification: OK}. The second run also shows \texttt{depth=3} (ISRG Root X1)
            in its verify lines.

        \subsection{Exercise 2: What does the \texttt{verify error} mean?}

            \begin{codeblock}[]{text}{First command (no trust store)}
                depth=2 C=US, O=ISRG, CN=Root YR
                verify error:num=20:unable to get local issuer certificate
                ...
                Verification error: unable to get local issuer certificate
            \end{codeblock}

            Error 20 means OpenSSL could not build a path from the server certificate up to a
            certificate it trusts. The last certificate in the chain was
            \texttt{Root YR}, and its issuer (ISRG Root X1) is not in any trust store the client
            knows about. The connection is still encrypted, but the client has no proof that it
            is talking to the real twitter.com. A man in the middle could present a self-made
            chain and the client could not tell the difference. This is why the error must not
            be ignored in practice.

        \subsection{Exercise 3: The server certificate}

            The server certificate was saved as \texttt{twitter\_com.cert} (the first
            \texttt{BEGIN/END CERTIFICATE} block) and displayed with
            \texttt{openssl x509 -in twitter\_com.cert -text}. The important fields are:

            \begin{codeblock}[]{text}{openssl x509 -in twitter\_com.cert -text (excerpt)}
                Version: 3 (0x2)
                Serial Number: 05:1c:d1:ae:fb:de:89:df:b5:9b:40:37:86:aa:c2:cf:0a:2f
                Signature Algorithm: sha256WithRSAEncryption
                Issuer: C=US, O=Let's Encrypt, CN=YR1
                Validity
                    Not Before: Sep 28 19:44:38 2026 GMT
                    Not After : Dec 27 19:44:37 2026 GMT
                Subject: CN=*.twitter.com
                Subject Public Key Info: rsaEncryption, 2048 bit, Exponent 65537
                X509v3 Key Usage: critical  Digital Signature, Key Encipherment
                X509v3 Extended Key Usage:  TLS Web Server Authentication
                X509v3 Basic Constraints: critical  CA:FALSE
                X509v3 Subject Key Identifier:   82:2A:D6:A3:...:1C:D3:4A
                X509v3 Authority Key Identifier: 1F:2F:35:BE:...:68:FB
                Authority Information Access: CA Issuers - URI:http://yr1.i.lencr.org/
                X509v3 Subject Alternative Name:
                    DNS:*.twitter.com, DNS:cdn.syndication.twitter.com, DNS:twitter.com
                X509v3 Certificate Policies: Policy: 2.23.140.1.2.1
                X509v3 CRL Distribution Points: URI:http://yr1.c.lencr.org/48.crl
                CT Precertificate SCTs: (2 timestamps)
                Signature Algorithm: sha256WithRSAEncryption  Signature Value: ...
            \end{codeblock}

            \begin{table}[h]
                \centering
                \small
                \begin{tabular}{p{0.28\textwidth}p{0.62\textwidth}}
                    \toprule
                    Field & Meaning \\
                    \midrule
                    Version & X.509 version (v3 allows extensions). \\
                    Serial Number & Unique number given by the issuer; used in revocation lists. \\
                    Signature Algorithm & Algorithm the issuer used to sign this certificate. \\
                    Issuer & The CA that signed the certificate (Let's Encrypt YR1). \\
                    Validity & Time window in which the certificate may be trusted (90 days here). \\
                    Subject & The identity being certified (\texttt{*.twitter.com}). \\
                    Subject Public Key Info & The subject's public key and its algorithm (RSA-2048). \\
                    Key Usage / Extended Key Usage & What the key may do: sign and encrypt keys, TLS server authentication. \\
                    Basic Constraints & \texttt{CA:FALSE}: this certificate cannot issue other certificates. \\
                    Subject/Authority Key Identifier & Hash of this key and of the issuer's key; helps find the issuer. \\
                    Authority Information Access & Where to download the issuer's certificate. \\
                    Subject Alternative Name & Hostnames the certificate is valid for. Modern clients check this, not the CN. \\
                    Certificate Policies & Policy OIDs; \texttt{2.23.140.1.2.1} is a domain-validated certificate. \\
                    CRL Distribution Points & Where to fetch the revocation list. \\
                    SCTs & Proof the certificate was logged in Certificate Transparency logs. \\
                    Signature Value & The issuer's digital signature over all fields above. \\
                    \bottomrule
                \end{tabular}
                \caption{Fields of an X.509 certificate.}
            \end{table}

        \subsection{Exercise 4: Is there an intermediate certificate?}

            Yes. The issuer of the server certificate is \texttt{CN=YR1, O=Let's Encrypt}, which is
            not a root. The AIA field gives its download location, so I downloaded it:

            \begin{codeblock}[]{bash}{Intermediate certificate}
                curl -o intermediate.cert http://yr1.i.lencr.org/
                openssl x509 -inform der -in intermediate.cert -text
            \end{codeblock}

            \begin{codeblock}[]{text}{Excerpt}
                Issuer: C=US, O=ISRG, CN=Root YR
                Validity: Sep  3 00:00:00 2025 GMT  to  Sep  2 23:59:59 2028 GMT
                Subject: C=US, O=Let's Encrypt, CN=YR1
                X509v3 Key Usage: critical  Digital Signature, Certificate Sign, CRL Sign
                X509v3 Basic Constraints: critical  CA:TRUE, pathlen:0
                Authority Information Access: CA Issuers - URI:http://yr.i.lencr.org/
            \end{codeblock}

            It is an intermediate CA (\texttt{CA:TRUE}, \texttt{pathlen:0}). Its purpose is to sign
            end-entity certificates on behalf of the root. The root private key can then stay
            offline and rarely used. If the intermediate is compromised, only that
            intermediate must be revoked, not the root. \texttt{pathlen:0} means it can sign
            leaf certificates but not further CAs. The file's SHA-256 fingerprint
            equals that of the second certificate sent by the server, so the server was already
            providing it.

        \subsection{Exercise 5: Is more than one organization involved?}

            Yes, there are separate certificate subjects. The chain
            presented by the server is:

            \begin{codeblock}[]{text}{Chain (from s\_client -showcerts)}
                0 s:CN=*.twitter.com            i:O=Let's Encrypt, CN=YR1
                1 s:O=Let's Encrypt, CN=YR1     i:O=ISRG, CN=Root YR
                2 s:O=ISRG, CN=Root YR          i:O=Internet Security Research Group, CN=ISRG Root X1
                (in the root store)
                3 s:O=Internet Security Research Group, CN=ISRG Root X1   (self-signed)
            \end{codeblock}

            The server certificate is signed by an intermediate named Let's Encrypt (the CA that
            serves the public), which is signed by Root YR, which is in turn signed by the
            self-signed root ISRG Root X1 operated by the Internet Security Research Group.
            These are different certificate subjects with different \texttt{O} values. Let's
            Encrypt is the service run by ISRG, so it is one organization family, but there are
            two levels of intermediates between the server and the root. Twitter itself is only
            the subject, not a CA. The reason for this design: the root key is kept offline for
            safety, and day-to-day signing is done by intermediates.

        \subsection{Exercise 6: Role of \texttt{ca-certificates.crt}}

            It is the \emph{trust store} (root store): a file of root CA certificates that the
            client trusts. During verification, OpenSSL builds a chain from the server
            certificate to a certificate in this file and checks every signature, validity
            period, and constraint along the way. A chain that ends at a certificate in this
            file is accepted. Everything else is rejected. Trust in all websites is therefore
            reduced to trust in the roots listed here.

        \subsection{Exercise 7: How many certificates are in the file?}

            \begin{codeblock}[]{bash}{Counting}
                grep -c "BEGIN CERTIFICATE" ca-certificates.crt
                openssl storeutl -noout -text -certs ca-certificates.crt | grep "Total found"
            \end{codeblock}

            \begin{codeblock}[]{text}{Output}
                121
                Total found: 121
            \end{codeblock}

            The Mozilla bundle I downloaded contains \textbf{121} root certificates. The count
            depends on the bundle version, so the course file may give a different number.

        \subsection{Exercise 8: A root certificate}

            I extracted \texttt{ISRG Root X1} (the root at the top of twitter's chain) from the
            bundle with a small Python script that splits the PEM blocks and checks each
            subject, and displayed it with \texttt{openssl x509 -in root\_isrg\_x1.pem -text}.

            \begin{table}[h]
                \centering
                \small
                \begin{tabular}{p{0.15\textwidth}p{0.24\textwidth}p{0.24\textwidth}p{0.24\textwidth}}
                    \toprule
                    & Twitter certificate & Intermediate (YR1) & Root (ISRG Root X1) \\
                    \midrule
                    Subject & CN=*.twitter.com & O=Let's Encrypt, CN=YR1 & O=Internet Security Research Group, CN=ISRG Root X1 \\
                    Issuer & O=Let's Encrypt, CN=YR1 & O=ISRG, CN=Root YR & \textbf{same as its subject} \\
                    Valid until & Dec 27 2026 & Sep 2 2028 & Jun 4 2035 \\
                    Basic Constraints & CA:FALSE & CA:TRUE, pathlen:0 & CA:TRUE \\
                    Key usage & Digital Signature, Key Encipherment & Digital Signature, Cert Sign, CRL Sign & Cert Sign, CRL Sign \\
                    \bottomrule
                \end{tabular}
                \caption{Comparison of the three certificates.}
            \end{table}

            Yes, the root has an Issuer field, and it is identical to its Subject. It is
            \emph{self-signed}: the root vouches for itself, so nobody above it can vouch for
            it. Its trust comes only from being included in the root store. The leaf's issuer
            equals the subject of the intermediate, and the intermediate's issuer
            names the next certificate up, which gives the chain of trust. The
            root has the longest validity (to 2035), the intermediate a few years,
            and the leaf only 90 days, since keys that are used more often are replaced sooner.

        \subsection{Exercise 9: DER to PEM}

            The intermediate downloaded from the AIA URL was in DER format (binary), so it was
            converted to PEM:

            \begin{codeblock}[]{bash}{Conversion}
                openssl x509 -inform der -in intermediate.cert -out intermediate.pem
            \end{codeblock}

    \section{Exercise 10: Certificate validation in Python}

        I used the supplied code (PyOpenSSL \texttt{X509Store} and \texttt{X509StoreContext}),
        modified to take the host name, fetch the chain from the server with
        \texttt{openssl s\_client -showcerts}, and treat the first certificate as the target and the
        others as intermediates. Each host was tested in two ways: with the root store
        only, and with the root store plus the intermediates.

        \begin{codeblock}[]{python}{verify\_chain.py (core)}
            def verify_chain_of_trust(cert_pem, trusted_cert_pems):
                certificate = crypto.load_certificate(crypto.FILETYPE_PEM, cert_pem)
                store = crypto.X509Store()
                for trusted_cert_pem in trusted_cert_pems:
                    trusted_cert = crypto.load_certificate(crypto.FILETYPE_PEM, trusted_cert_pem)
                    store.add_cert(trusted_cert)
                store_ctx = crypto.X509StoreContext(store, certificate)
                try:
                    store_ctx.verify_certificate()   # None means validated
                    return True, "ok"
                except crypto.X509StoreContextError as e:
                    return False, str(e)

            def verify(host, use_intermediate=True):
                chain = fetch_chain(host)               # s_client -showcerts
                cert, int_certs = chain[0], chain[1:]
                trusted = [str(p) for p in pem.parse_file("./ca-certificates.crt")]
                if use_intermediate:
                    trusted.extend(int_certs)
                return verify_chain_of_trust(cert, trusted)
        \end{codeblock}

    \pagebreak

        \begin{codeblock}[]{text}{Output of python verify\_chain.py}
            twitter.com       [root store only ] -> FAILED (unable to get local issuer certificate)
            twitter.com       [with intermediate] -> Certificate verified
            google.com        [root store only ] -> FAILED (unable to get local issuer certificate)
            google.com        [with intermediate] -> Certificate verified
            www.chula.ac.th   [root store only ] -> FAILED (unable to get local issuer certificate)
            www.chula.ac.th   [with intermediate] -> Certificate verified
            classdeedee.cloud.cp.eng.chula.ac.th [root store only ] -> FAILED (unable to get local issuer certificate)
            classdeedee.cloud.cp.eng.chula.ac.th [with intermediate] -> Certificate verified
        \end{codeblock}

        \begin{table}[h]
            \centering
            \small
            \begin{tabular}{p{0.27\textwidth}p{0.33\textwidth}p{0.25\textwidth}}
                \toprule
                Host & Chain sent by server (leaf issuer $\rightarrow$ \ldots) & Trust anchor in bundle \\
                \midrule
                twitter.com & Let's Encrypt YR1 $\rightarrow$ ISRG Root YR & ISRG Root X1 \\
                google.com & Google Trust Services WR2 $\rightarrow$ GTS Root R1 & GlobalSign Root CA \\
                www.chula.ac.th & Sectigo DV R36 $\rightarrow$ Sectigo Root R46 & USERTrust RSA Certification Authority \\
                classdeedee.cloud.cp.eng.\allowbreak chula.ac.th & Let's Encrypt YR2 $\rightarrow$ ISRG Root YR & ISRG Root X1 \\
                \bottomrule
            \end{tabular}
            \caption{Chain of each host and the trust anchor that validates it.}
        \end{table}

        All four certificates were verified when the intermediates were supplied. Without the
        intermediates, validation fails for every host, because the root store contains only
        roots and the issuer of the leaf is an intermediate that the store does not know.
        This is why servers must send the intermediate certificates along with their own.
        Note that the classdeedee server returns the certificate
        \texttt{CN=cloud.cp.eng.chula.ac.th}, a certificate of the parent domain. The script checks the
        chain only: it does not check that the name matches the host or that the certificate is
        not revoked.

    \pagebreak

    \section{Exercise 11: Class 1 and Class 3 certificates}

        The \emph{class} describes how strongly the CA checks the identity of the applicant
        before it signs (the assurance level of the certificate policy).

        \textbf{Class 1} (low assurance): only a minimal check, usually that the applicant
        controls an email address. No identity documents are required. Typical uses:
        personal email signing and encryption (S/MIME), and low-risk logins where the user
        only needs to be recognised, not identified. It is cheap and quick but cannot prove
        who the person or organisation really is, so it is unsuitable for banking or commerce.

        \textbf{Class 3} (high assurance): the CA verifies the real identity of the person or
        legal entity with documents, often in person or with a notary, and checks the
        applicant's authority. Typical uses: certificates for commercial web servers, code
        signing, and certificates for organisations and CAs. It is more expensive and slower
        but allows relying parties to trust that the named organisation is real.

        A Class~1 certificate has uses that Class~3 does not need: free personal email
        protection where convenience matters and identity does not. A Class~3 certificate
        is needed where a wrong identity causes damage (e-commerce, signed software), and a
        Class~1 certificate would not be accepted by relying parties there. Today these
        classes have largely been replaced by validation types (DV, OV, EV), which follow the
        same idea.

    \pagebreak

    \section{Exercise 12: A compromised root CA}

        \subsection{12a: Further attacks}

            If a root CA in the client's root store is controlled by an attacker, the attacker
            can sign a certificate for \emph{any} name, and every client that trusts that root
            will accept it as valid. Possible attacks:

            \begin{itemize}
                \item \textbf{Man-in-the-middle on HTTPS:} issue a certificate for
                    \texttt{bank.example}, then intercept, read, and modify encrypted traffic with no
                    browser warning.
                \item \textbf{Impersonation and phishing:} fake sites with a valid padlock.
                \item \textbf{Malware and updates:} sign malicious code or updates that seem
                    to come from a trusted vendor, if the root also covers code signing.
                \item \textbf{Forged email, documents, and client certificates} (S/MIME,
                    document signing, certificate-based login).
                \item \textbf{Targeted or mass interception} (for example by a hostile
                    network), which can go unnoticed for months because the forged certificates
                    look valid.
            \end{itemize}

            \begin{figure}[h]
                \centering
                \begin{tikzpicture}[
                    node distance=2.2cm and 2.4cm,
                    box/.style={draw, rounded corners, minimum width=2.4cm, minimum height=1cm, align=center, fill=blue!6},
                    bad/.style={draw, rounded corners, minimum width=2.4cm, minimum height=1cm, align=center, fill=red!12},
                    arr/.style={-{Stealth}, thick}
                ]
                    \node[box] (client) {Victim\\browser};
                    \node[bad, right=of client] (mitm) {Attacker\\(man in the middle)};
                    \node[box, right=of mitm] (server) {Real server\\bank.example};
                    \node[bad, above=of mitm] (root) {Compromised\\root CA};
                    \node[box, above=of client] (store) {Root store\\(trusts root)};

                    \draw[arr, <->] (client) -- node[above, font=\scriptsize, align=center] {TLS\\(forged cert)} (mitm);
                    \draw[arr, <->] (mitm) -- node[above, font=\scriptsize, align=center] {TLS\\(real cert)} (server);
                    \draw[arr, dashed] (root) -- node[right, font=\scriptsize, align=left] {signs fake cert\\for bank.example} (mitm);
                    \draw[arr, dashed] (root) -- node[above, font=\scriptsize] {trusted by} (store);
                \end{tikzpicture}
                \caption{Attack setup: the attacker uses a certificate signed by the compromised root and relays traffic to the real server.}
            \end{figure}

        \subsection{12b: Can we rely on CRLs or OCSP?}

            \textbf{No, not reliably.}
            \begin{itemize}
                \item Revocation information is signed by the CA. If the root key or its
                    infrastructure is controlled by the attacker, the attacker can also
                    sign CRLs and OCSP responses, marking the forged certificates as
                    ``good'', or simply never publish a revocation.
                \item A root certificate is self-signed and has no issuer above it that could
                    revoke it. Clients do not check CRL/OCSP for the trust anchor. The only
                    fix is out-of-band: the browser or OS vendor removes the root from the
                    root store with an update.
                \item In practice, many clients \emph{soft-fail} (accept the certificate when
                    the OCSP server is unreachable), so an attacker in the middle can simply
                    block OCSP requests. CRLs also have a delay between revocation and
                    distribution.
                \item Since the compromise is unnoticed for months, nothing would be revoked
                    in the meantime anyway.
            \end{itemize}

            Additional defences are needed: removing the root from trust stores, Certificate
            Transparency monitoring (forged certificates become visible in logs), certificate
            or key pinning, HSTS, and limiting what each root may issue (name constraints).

\end{document}
